\documentclass[11pt,a4paper]{article}
\usepackage{DDTA_METHODOLOGY_GUIDE_STYLE_R1}
\usepackage{paracol}

\hypersetup{
  pdftitle={DDTA DermaTriage - Documentation and Base Analysis Case Study},
  pdfauthor={Documentation-Driven Threat Analysis research project}
}

% -----------------------------------------------------------------------------
% PAGE-INTEGRITY HASHES
% Generated over the canonical LaTeX payload between PAGE-CONTENT markers.
% Hash lines themselves are excluded. On the integrity-index page, hashes of
% preceding pages are resolved before hashing, while the page's own hash is
% represented by the canonical token <SELF> to avoid self-reference.
% -----------------------------------------------------------------------------
\newcommand{\PageMDfive}[1]{\csname PageMDfive#1\endcsname}
%<PAGE-HASH-DEFS>
\expandafter\def\csname PageMDfive1\endcsname{ccabe4bed7cca4d2f68a2005e42f36e3}
\expandafter\def\csname PageMDfive2\endcsname{0fb76dc0491da2793823b20571e5268b}
\expandafter\def\csname PageMDfive3\endcsname{54275f1583e53f3d9e98664420b2b6c6}
%</PAGE-HASH-DEFS>

\newcommand{\PageHashFooter}[1]{%
  \fancyfoot[R]{\sffamily\scriptsize\color{DDTAGray}MD5 contenuto: \texttt{#1}}%
}

\newcommand{\IntegrityRow}[3]{#1 & #2 & {\footnotesize\ttfamily #3} \\
}

\begin{document}
\selectlanguage{italian}
\fancyhead[L]{\sffamily\small\color{DDTAGray}DDTA - DermaTriage Case Study}
\fancyhead[R]{\sffamily\small\color{DDTABlue}Documentation + Base Analysis}
\fancyfoot[L]{\sffamily\scriptsize\color{DDTAGray}Documentation 2/3 - Base Analysis 1/3}

%<PAGE-DESC:1>Frontespizio e struttura del case study
%<PAGE-CONTENT:1>
\begin{titlepage}
\thispagestyle{empty}
\vspace*{10mm}
{\sffamily\bfseries\color{DDTANavy}\fontsize{15}{18}\selectfont Documentation-Driven Threat Analysis}\par
\vspace{5mm}
{\sffamily\bfseries\color{DDTANavy}\fontsize{29}{33}\selectfont DermaTriage Case Study}\par
\vspace{2mm}
{\sffamily\color{DDTABlue}\fontsize{18}{22}\selectfont Documentation + Base Analysis}\par

\vspace{12mm}
\begin{tcolorbox}[enhanced,arc=3pt,boxrule=0pt,colback=DDTANavy,left=13pt,right=13pt,top=12pt,bottom=12pt]
{\color{white}\sffamily\bfseries\large Ricostruzione DDTA del progetto DermaTriage}\par
\vspace{2mm}
{\color{white!88}\small La documentazione di progetto occupa la parte principale della pagina. La Base Analysis viene riportata in parallelo nella colonna laterale soltanto quando esiste un significato documentale da analizzare.}
\end{tcolorbox}

\vfill
\begin{ddtacore}[title={Struttura del case study}]
Il documento segue la progressione DDTA dal contesto generale del progetto agli elementi documentali successivi. Ogni elemento viene presentato nella colonna documentale; la colonna Base Analysis mostra esclusivamente l'analisi derivabile da quel significato, senza introdurre fatti di progetto aggiuntivi.
\end{ddtacore}

\vspace{5mm}
{\sffamily\scriptsize\color{DDTAGray}MD5 contenuto pagina 1: \texttt{\PageMDfive{1}}}\par
\end{titlepage}
%</PAGE-CONTENT:1>

%<PAGE-DESC:2>Contesto generale del progetto - Project Problem Framing
%<PAGE-CONTENT:2>
\PageHashFooter{\PageMDfive{2}}
\section{Contesto generale del progetto}

\setlength{\columnsep}{7mm}
\setlength{\columnseprule}{0.35pt}
\columnratio{0.67}
\begin{paracol}{2}

% LEFT 2/3 - DOCUMENTATION
{\sffamily\large\bfseries\color{DDTANavy}Documentazione DDTA}\par
\vspace{0.5em}

\begin{ddtacore}[title={Project Problem Framing}]
Il problema che DermaTriage deve affrontare e' il supporto al triage precoce di casi dermatologici relativi a lesioni potenzialmente oncologiche. Il progetto parte dalle informazioni gia' disponibili sul caso per determinarne l'urgenza e una priorita' operativa di presa in carico, con l'obiettivo di favorire un instradamento specialistico tempestivo.
\end{ddtacore}

\switchcolumn

% RIGHT 1/3 - BASE ANALYSIS ONLY
{\sffamily\large\bfseries\color{DDTABlue}Base Analysis}\par
\vspace{0.5em}

\begin{ddtaworking}[title={Base Analysis}]
\small
Il contesto generale definisce il problema affrontato dal progetto, ma non introduce ancora un elemento documentale strutturato da decomporre in referenti e proposizioni. La Base Analysis viene quindi applicata a partire dagli elementi DDTA successivi, quando il significato governato rende identificabili responsabilita', relazioni, condizioni o risultati analizzabili.
\end{ddtaworking}

\end{paracol}
%</PAGE-CONTENT:2>

\clearpage
%<PAGE-DESC:3>Indice di integrita' delle pagine
%<PAGE-CONTENT:3>
\PageHashFooter{\PageMDfive{3}}
\section{Indice di integrita' delle pagine}

\begin{ddtacore}[title={Regola di verifica}]
Ogni pagina possiede un MD5 del proprio contenuto canonico. Il calcolo esclude il valore MD5 stampato sulla pagina, gli header/footer e i metadati di impaginazione esterni al relativo blocco di contenuto. Per la pagina dell'indice, gli hash delle pagine precedenti fanno parte del contenuto controllato; il solo hash della pagina indice viene sostituito dal token canonico \texttt{<SELF>} durante il calcolo, evitando una dipendenza circolare.
\end{ddtacore}

\vspace{3mm}
\begin{tabularx}{\textwidth}{L{12mm}L{72mm}Y}
\toprule
\textbf{Pag.} & \textbf{Contenuto} & \textbf{MD5 contenuto} \\
\midrule
%<PAGE-INTEGRITY-ROWS>
\IntegrityRow{1}{Frontespizio e struttura del case study}{\PageMDfive{1}}
\IntegrityRow{2}{Contesto generale del progetto - Project Problem Framing}{\PageMDfive{2}}
\IntegrityRow{3}{Indice di integrita' delle pagine}{\PageMDfive{3}}
%</PAGE-INTEGRITY-ROWS>
\bottomrule
\end{tabularx}

\vspace{4mm}
{\small\color{DDTAGray}
Quando una pagina viene modificata intenzionalmente, cambia il suo MD5. Le pagine non interessate dal passaggio devono mantenere lo stesso valore; una variazione inattesa segnala una modifica da riesaminare prima di proseguire.}
%</PAGE-CONTENT:3>

\end{document}
